% Generated draft: requires mathematical and source review.
\documentclass[11pt]{article}
\usepackage[margin=25mm]{geometry}
\usepackage{fontspec}
\setmainfont{Latin Modern Roman}
\usepackage{amsmath,amssymb,mathtools,unicode-math}
\setmathfont{Latin Modern Math}
\usepackage{enumitem,needspace,xurl,hyperref}
\setlength{\parindent}{0pt}
\setlength{\parskip}{5pt}
\setlength{\emergencystretch}{3em}
\newcommand{\N}{\mathbb{N}}
\newcommand{\Z}{\mathbb{Z}}
\newcommand{\Q}{\mathbb{Q}}
\newcommand{\R}{\mathbb{R}}
\newcommand{\C}{\mathbb{C}}
\newcommand{\F}{\mathbb{F}}
\newcommand{\A}{\mathbb{A}}
\newcommand{\PP}{\mathbb{P}}
\newcommand{\E}{\mathbb{E}}
\newcommand{\eps}{\varepsilon}
\newcommand{\dd}{\,\mathrm{d}}
\begin{document}
\section*{UnsolvedMath ID: 800473}
\textbf{Source problem code:} catalogue entry
\subsection{Definitions and mathematical statement}
The source poses an exploratory construction problem rather than a specified existence theorem. Its motivating discrete system consists of three point particles, called beads, with masses \(m_1,m_2,m_3\), moving on a frictionless ring and undergoing elastic collisions. Write \(M=m_1+m_2+m_3\). The momentum and kinetic energy used in this motivation are
\[
P=\sum_{i=1}^{3}m_i v_i,\qquad E=\frac12\sum_{i=1}^{3}m_i v_i^2,
\]
where \(v_i\) is the velocity of bead \(i\); negative masses make the energy potentially indefinite.

For positive masses, the source states that conservation of momentum and energy, followed by quotienting out rotational symmetry, identifies the motion with billiard motion in an acute triangle whose angles satisfy
\[
\tan\alpha_i=\sqrt{\frac{m_iM}{m_jm_k}},\qquad (i,j,k)\in\{(1,2,3),(2,3,1),(3,1,2)\}.
\]
It describes the limit \(m_3\to+\infty\) as a right triangular billiard with \(\alpha_3=\pi/2\), equivalent to two beads on a segment between two walls. The stated correspondence motivates questions about periodic trajectories.

The source also permits negative masses subject to
\[
Mm_1m_2m_3>0.
\]
Under this restriction, it reports that masses of the same sign correspond to an acute triangle, whereas mixed signs correspond to an obtuse triangle. The positive square-root tangent formula above is stated here only for positive masses: an obtuse-angle convention for the signed-mass correspondence is not supplied. An alternative viewpoint in the source keeps the triangle fixed and changes its metric by rescaling; the negative-mass setting is described there using pseudo-Euclidean metrics. A pseudo-Euclidean metric means a constant nondegenerate indefinite symmetric bilinear form; a pseudo-Riemannian metric is its smoothly varying analogue. The source does not specify the reduced metric or its signature for each admissible sign pattern.

\textbf{Requested construction.} Construct a continuous version of the interaction theory of particles of positive and negative mass, allowing a mass density that takes negative values, and explore its relation to pseudo-Euclidean geometry and, where appropriate, pseudo-Riemannian geometry. Formulate the corresponding Euler-like fluid evolution equations and investigate whether they exhibit interesting phenomena associated with negative density. The continuum model should explain its connection to the motivating signed-mass particle dynamics.

Here a density that takes negative values is the intended interpretation of ``negative density,'' but the source does not decide between an everywhere negative density and a density with both positive and negative regions. It specifies no spatial domain, dimension, interaction law, continuum limiting procedure, constitutive assumptions, initial or boundary conditions, or solution concept. Consequently, no particular system of partial differential equations or quantified existence assertion can be extracted without additional choices.

The motivating discussion also reports Artigue's study of positive and negative masses and describes a small negative-mass particle, called a ``graviton,'' that pulls positive masses together through elastic collisions. Its suggestion that this might account for gravity is speculative motivation, not an established assertion or a required conclusion of this problem.
\subsection{Short English statement}
Construct a continuum theory for interacting positive and negative masses, identify suitable Euler-like fluid equations, and study their relation to pseudo-Euclidean geometry. The source motivates this through elastic bead collisions and triangular billiards, but leaves the fluid model, meaning of negative density, and precise mathematical targets unspecified.
\subsection{Sources}
\begin{itemize}
\item Local UnsolvedMath registry, permanent ID \(800473\), exact source problem code \texttt{catalogue entry}. The supplied record labels the problem \emph{open}; this is a supplied registry classification, not an independently verified mathematical or literature assessment.
\item Albach, Barraud et al., \emph{Open Problems in Billiards and Quantitative Symplectic Geometry} (2026), \url{https://arxiv.org/abs/2602.12896}. The supplied attribution identifies Boris Khesin as contributor and the source item as ``Boris Khesin --- Problem 2.'' The extracted statement itself is headed ``3. Is there a fluid of negative density?'' Supplied access date: 2026-10-16; supplied extraction method: \texttt{arxiv\_latex}. These provenance details have not been independently verified.
\item Reference label [GlashowMittag]: cited in the supplied statement for the negative-mass mechanical correspondence with obtuse triangular billiards. No full bibliographic entry or separate URL is supplied.
\item Reference label [Artigue]: cited for a study of positive and negative masses, the negative-mass ``graviton'' collision mechanism, and the associated theory and applications. No full bibliographic entry or separate URL is supplied.
\item Reference label [Wiki]: cited in connection with the speculative gravity motivation. No title, bibliographic entry, or URL is supplied.
\item Reference label [KhTab]: cited for billiards in pseudo-Euclidean and more general pseudo-Riemannian metrics. No full bibliographic entry or separate URL is supplied.
\item The supplied acknowledgments state that Boris Khesin thanks the organizers and participants of the workshop on \emph{Billiards and Quantitative Symplectic Geometry} at Heidelberg University and acknowledges partial support from an NSERC Discovery Grant.
\end{itemize}
\end{document}
